
\section{Related Work}

We review prior work on design patterns of composite visualizations,
predictability and flexibility in visualization authoring, and interactive
mechanisms for composing visual structures.


\subsection{Design Patterns of Composite Visualization}



Multiple-view and composite visualizations combine distinct visual
representations to support comparison, coordination, and analysis of complex
data. Prior work has established guidelines for coordinated multiple views and
identified recurring comparison strategies such as juxtaposition,
superposition, and explicit encoding~\cite{GuidelinesMultipleViews,
CoordinatedMultipleViews,VisualComparison}. Related design spaces further
examine dashboard layouts, domain-specific multi-view structures, and
multifaceted visualizations~\cite{DashboardDesignSpace,
MultiViewDesignPatterns,MultiFacetedGraph}.

More directly, composite visualization taxonomies characterize recurring
spatial, semantic, and data relationships among views. Early work identifies
patterns including juxtaposition, integration, superimposition, overloading,
and nesting~\cite{javed2012exploring}. Subsequent studies distinguish view
composition from spatial configuration~\cite{Composite}, derive finer
taxonomies based on spatial and semantic relationships~\cite{deng2023survey},
and index composite designs by their visual structures, data characteristics,
and analytical tasks~\cite{vaid}. Together, these studies provide a vocabulary
for describing how completed composite visualizations are organized.

These models, however, do not determine the authoring units through which such
structures should be constructed. {\system} builds on these taxonomies by
using spatial relationships as references for direct manipulation, while
grounding the decomposition into chart-level blocks in how visualization
experts perceive composite structures.


\subsection{Predictability and Flexibility in Visualization Authoring}
\label{sec:related-authoring-granularity}

Visualization authoring approaches make different trade-offs between
predictability and flexibility, depending on how explicitly visual structures
are presented to authors and how freely they can be adapted.

Fine-grained graphical authoring systems~\cite{D3,Protovis,Lyra,Lyra2,
DataIllustrator,Charticulator,iVisDesigner,DataDrivenGuides,StructGraphics,
Hanpuku,iVoLVER,DataInk,DataBreeze,InChorus,Liger,
VisualizationbyDemonstration,CreatingChartsbyDemonstrationGold} provide
further flexibility by exposing graphical and structural components below the
chart level. Data Illustrator~\cite{DataIllustrator} and
Charticulator~\cite{Charticulator} support highly customized structures, but
place greater construction and coordination demands on authors.


Explicit visual-template systems~\cite{ManyEyes,RAWGraphs} present
recognizable chart forms with bounded data roles, allowing authors to
anticipate the resulting visualization before configuration.
Many Eyes~\cite{ManyEyes} and RAWGraphs~\cite{RAWGraphs} exemplify this
chart-chooser paradigm, where the visual form is explicit from the outset.
Such templates minimize structural reasoning and provide highly predictable
outcomes, but their expressiveness is limited by the available template
collection.

Generalized grammar- and schema-based systems~\cite{Tableau, Vega,VegaLite,
ggplot2LayeredGrammarofGraphics,Polaris,Tableau,ApacheECharts,Altair,Ivy}
increase flexibility by abstracting across chart classes through shared marks,
encoding channels, structural roles, or parameters. Vega-Lite~\cite{VegaLite}
systematically combines these abstractions to cover a broad design space,
while Ivy~\cite{Ivy} generalizes examples through parameterized templates.
Such abstractions reduce the need for a separate template for every chart
variant, but require authors to reason about the underlying grammar and
compatible design choices rather than selecting an immediately explicit
visual outcome.



{\system} prioritizes the visual clarity and predictability of explicit
templates by exposing recognizable chart blocks as authoring units. Rather
than replacing these units with a more abstract grammar, related blocks are
organized into interoperable chart families, whose boundaries are grounded in
how experts perceive constituent units of composite visualizations. This
provides controlled flexibility across related structures while preserving
the explicit identity of each block.

\subsection{Interactive Composition Mechanisms}
\label{sec:related-composition-systems}

Given explicit authoring units, existing systems differ in how authors express
and manipulate composition relationships among them.

Operator-based composition systems~\cite{VegaLite,Vega,
ggplot2LayeredGrammarofGraphics,Altair,ApacheECharts,Atlas,ATOM,
ComputableViz,GoFish,Gosling,GenomeSpy,GoTree,Cicero,Nebula,Mosaic,Encodable}
encode composition through operators and specification structures.
Vega-Lite~\cite{VegaLite} provides operators for layering,
concatenation, and faceting, while GoFish~\cite{GoFish} supports recursive
graphical composition. These mechanisms provide composition semantics
and broad coverage, but authors express relationships through specification
abstractions rather than directly manipulating chart units.


Visual assembly systems~\cite{SnapTogetherVisualization,Improvise,VisMashup,
VisComposer,SemanticSnapping,Medley,MultiVision,DrillboardsDrillVis,
QuickDashboard,QualDash,Pilingjs,GEViTRec} make views or components
manipulable during composition. Semantic Snapping~\cite{
SemanticSnapping} guides the combination of compatible views, while
Medley~\cite{Medley} supports assembly of multi-view designs.
These approaches make constituent views and their spatial relationships clear,
but primarily focus on arrangement and coordination among complete views.

Structure-aware composition systems~\cite{Visception,CompositingVis} extend
direct composition to hierarchical or internal visual structures.
Visception~\cite{Visception} supports chart-in-chart composition through
hierarchical mappings, while CompositingVis~\cite{CompositingVis} maps spatial
manipulations to different composite-view relationships. These systems make
more complex composition structures directly authorable, although their
interactions are designed around particular hierarchical or spatial
composition settings.

{\system} combines direct manipulation with explicit composition semantics at
the chart-block level. Dragging one block relative to another uses
layout-referenced targets to express concatenation, layering, faceting, or
nesting, including targets within another chart. This provides a consistent
interaction across composition forms while keeping constituent blocks and
their relationships explicit and editable.
